\subsection{Probing Physical Information in Latent Representations}
\label{sec:mid_interventions}

\begin{figure}[H]
\centering
\includegraphics[width=\linewidth]{figures/latent_analysis.pdf}
\caption{\textbf{Physical information in latent representations.}
(a) Token representations in a shared principal component projection, colored by reasoning role.
(b) Row-normalized confusion matrix of the six-domain linear probe.
(c) Accuracy of predicting both outcomes in each same-image intervention pair; error bars show 95\% confidence intervals from resampling scene families.
(d) Mean physical question-answering accuracy before and after intermediate-token perturbations.}
\label{fig:latent_analysis}
\label{fig:mid_interventions}
\label{fig:latent_structure}
\end{figure}

\paragraph{Token roles remain distinct after compression.}
The four reasoning roles occupy distinct regions in a shared principal component projection (Figure~\ref{fig:latent_analysis}a), showing role-dependent structure after compression.

\paragraph{Physical domains are linearly decodable.}
We freeze PhysFlow and fit a class-balanced linear ridge classifier to predict PhysTrace's six physical domains from the concatenation of two textual and four intermediate visual tokens.
Four-fold cross-validation on 212 validation videos selects regularization; evaluation on 212 independent test videos yields \textbf{86.91\% balanced accuracy}, versus 16.67\% for uniform random guessing.
Balanced accuracy averages recall across the six domains (Figure~\ref{fig:latent_analysis}b).

\paragraph{Visual tokens distinguish intervention outcomes.}
We test outcome prediction with synthetic scenes covering 11 exploratory tasks.
Each pair keeps the initial image fixed and uses two instructions that select different objects or settings, producing opposite binary outcomes under idealized physical rules.
For example, removing the load-bearing support makes a ball fall, whereas removing the other support leaves it supported.
We fit separate linear ridge probes to the four frozen intermediate visual tokens, using disjoint training, validation, and test scene families (128/32/64 per task).
Each family contains two related source images; validation selects mean pooling or concatenation and ridge regularization.
A pair counts as correct only when the probe predicts both outcomes correctly.

Figure~\ref{fig:latent_analysis}c highlights seven high-scoring tasks selected after exploratory evaluation, spanning rigid-body, fluid, state-change, and deformation phenomena, with paired accuracy from \textbf{91.41\% to 100\%}.
The remaining four exploratory tasks yield 28.91\% to 68.75\%.
Each displayed task uses 128 test pairs, except phase change, for which the fixed probe correctly predicts all 256 pairs from 128 new confirmation families.
Randomly assigning the two opposite labels within a pair gives the 50\% reference; independent binary guesses give 25\%.

\paragraph{Token ordering affects physical editing.}
We evaluate 64 PICABench-Superficial inputs~\citep{pu2025picabench}, eight per category, while fixing the textual and final visual tokens.
After planning, we mask all four intermediate visual tokens, reverse their contents, or shuffle them as $(3,1,4,2)$.
Qwen2.5-VL-72B-Instruct~\citep{bai2025qwen25vl} answers physical validation questions on the edited images; we average accuracy within each input and then across inputs.
Reversing or shuffling reduces mean accuracy from 57.36\% to 53.94\%, a \textbf{3.42 percentage-point decrease} (Figure~\ref{fig:latent_analysis}d).

Together, the analyses establish physical-information accessibility and ordering sensitivity.
Input-based and textual probes also solve many tasks with visible material or temperature cues; token-to-frame temporal alignment remains unresolved.
